\chapter{Complexity — Layout}
%%% generated by the blueprint pipeline from TCSlib.Complexity.SpaceComplexity.Machines.Layout
\section{Overview}

Declarations of the Lean module \texttt{TCSlib.Complexity.SpaceComplexity.Machines.Layout}.

\section{Declarations}

\begin{definition}[Boolean segment]
\lean{Complexity.LogProg.Seg}
\label{Complexity.LogProg.Seg}
\leanok
A segment is a pair consisting of a finite Boolean word and a Boolean flag.
\end{definition}

\begin{definition}[Rendering a segment]
\lean{Complexity.LogProg.render}
\label{Complexity.LogProg.render}
\leanok
\uses{Complexity.LogProg.Seg, Turing.dbl}
For a segment consisting of a Boolean word $w$ and a flag $b$, its rendering is the
bitwise doubling of $w$ when $b$ is true, and is $w$ itself when $b$ is false.
\end{definition}

\begin{definition}[Rendered segment length]
\lean{Complexity.LogProg.rlen}
\label{Complexity.LogProg.rlen}
\leanok
\uses{Complexity.LogProg.Seg}
For a segment $(w,b)$, define its rendered length to be $2|w|$ if $b$ is true and $|w|$
otherwise, where $|w|$ is the length of the Boolean word $w$.
\end{definition}

\begin{lemma}[Length of a doubled segment]
\lean{Complexity.LogProg.rlen_true}
\label{Complexity.LogProg.rlen_true}
\leanok
\uses{Complexity.LogProg.rlen}
\difficulty{1}
For every Boolean word $w$, the length assigned to the segment $(w,\mathrm{true})$ is
$2|w|$.
\end{lemma}

\begin{lemma}[Length of an undoubled segment]
\lean{Complexity.LogProg.rlen_false}
\label{Complexity.LogProg.rlen_false}
\leanok
\uses{Complexity.LogProg.rlen}
\difficulty{1}
For every Boolean word $w$, the length assigned to the segment $(w,\mathrm{false})$ is
$|w|$.
\end{lemma}

\begin{lemma}[Length of a rendered segment]
\lean{Complexity.LogProg.length_render}
\label{Complexity.LogProg.length_render}
\leanok
\uses{Complexity.LogProg.Seg, Complexity.LogProg.render, Complexity.LogProg.rlen}
\difficulty{3}
For any segment consisting of a bit string $w$ and a flag indicating whether it is
doubled, the length of its rendered bit string equals $2|w|$ when the flag is true and
$|w|$ when it is false.
\end{lemma}

\begin{definition}[Encoding a list of segments]
\lean{Complexity.LogProg.vword}
\label{Complexity.LogProg.vword}
\leanok
\uses{Complexity.LogProg.Seg, Complexity.LogProg.render}
Define the encoding of a list of segments recursively: the empty list encodes as the
empty word, and a one-segment list encodes as the rendering of that segment. For a list
with at least two segments, encode the first segment, append the separator $01$, and
then encode the remaining list.
\end{definition}

\begin{definition}[Segment-list offset]
\lean{Complexity.LogProg.off}
\label{Complexity.LogProg.off}
\leanok
\uses{Complexity.LogProg.Seg, Complexity.LogProg.rlen}
Define the offset of a list of segments at a natural number $i$ recursively: it is $0$
when $i=0$ or when the list is empty; for a nonempty list and positive $i$, it is the
rendered length of the first segment plus $2$ plus the offset of the remaining list at
$i-1$.
\end{definition}

\begin{definition}[Segment lookup with a default]
\lean{Complexity.LogProg.seg}
\label{Complexity.LogProg.seg}
\leanok
\uses{Complexity.LogProg.Seg}
A segment is a pair consisting of a finite Boolean list and a Boolean flag. Given a list
of segments and an index, the selected segment is the entry at that index when it
exists, and otherwise is the pair consisting of the empty list and `false`.
\end{definition}

\begin{lemma}[Segment lookup at the list head]
\lean{Complexity.LogProg.seg_cons_zero}
\label{Complexity.LogProg.seg_cons_zero}
\leanok
\uses{Complexity.LogProg.Seg, Complexity.LogProg.seg}
\difficulty{1}
Let a segment be a pair consisting of a Boolean list and a Boolean. For any segment $s$
and list of segments $R$, looking up position zero in $s :: R$ returns $s$.
\end{lemma}

\begin{lemma}[Segment lookup after a list head]
\lean{Complexity.LogProg.seg_cons_succ}
\label{Complexity.LogProg.seg_cons_succ}
\leanok
\uses{Complexity.LogProg.Seg, Complexity.LogProg.seg}
\difficulty{1}
Let a segment be a pair consisting of a Boolean list and a Boolean. For any segment $s$,
list of segments $R$, and natural number $i$, looking up position $i+1$ in the list $s
:: R$ returns the same segment as looking up position $i$ in $R$.
\end{lemma}

\begin{lemma}[Successive segment offsets]
\lean{Complexity.LogProg.off_succ}
\label{Complexity.LogProg.off_succ}
\leanok
\uses{Complexity.LogProg.Seg, Complexity.LogProg.off, Complexity.LogProg.rlen, Complexity.LogProg.seg, Complexity.LogProg.seg_cons_succ}
\difficulty{4}
Let $segs$ be a finite list of segments and let $i$ satisfy $i+1<|segs|$. The offset at
index $i+1$ equals the offset at index $i$, plus the rendered length of segment $i$,
plus the two-bit separator length.
\end{lemma}

\begin{lemma}[Length of a virtual word]
\lean{Complexity.LogProg.length_vword}
\label{Complexity.LogProg.length_vword}
\leanok
\uses{Complexity.LogProg.Seg, Complexity.LogProg.length_render, Complexity.LogProg.off, Complexity.LogProg.rlen, Complexity.LogProg.seg, Complexity.LogProg.seg_cons_succ, Complexity.LogProg.vword}
\difficulty{5}
Let a nonempty finite list of segments be given. The length of the virtual word formed
by rendering the segments in order and placing the separator $01$ between consecutive
segments equals the sum of the rendered lengths of all segments except the last, plus
two for each such separator, plus the rendered length of the last segment.
\end{lemma}

\begin{lemma}[Segment Entries in the Virtual Word]
\lean{Complexity.LogProg.getElem?_vword_seg}
\leanok
\uses{Complexity.LogProg.Seg, Complexity.LogProg.off, Complexity.LogProg.render, Complexity.LogProg.rlen, Complexity.LogProg.seg, Complexity.LogProg.seg_cons_succ, Complexity.LogProg.seg_cons_zero, Complexity.LogProg.vword}
\difficulty{5}
Let a finite list of segments be given, and let $i$ be a valid segment index. If $k$ is
smaller than the rendered length of the segment at index $i$, then the possibly
undefined entry at index $\operatorname{off}(i)+k$ in the virtual word formed by
rendering the segments with separators equals the entry at index $k$ in the rendering of
segment $i$.
\end{lemma}

\begin{lemma}[Separators in an encoded bit word]
\lean{Complexity.LogProg.getElem?_vword_sep}
\leanok
\uses{Complexity.LogProg.Seg, Complexity.LogProg.off, Complexity.LogProg.render, Complexity.LogProg.rlen, Complexity.LogProg.seg, Complexity.LogProg.seg_cons_succ, Complexity.LogProg.seg_cons_zero, Complexity.LogProg.vword}
\difficulty{5}
Let $(s_i,\beta_i)$ be the $i$th segment of a finite list of Boolean segments, using the
default empty segment if the index is out of range. If $i+1$ is less than the number of
segments, then in the concatenated word that renders each segment and places the
separator $01$ between consecutive segments, the bit immediately after the rendered
$i$th segment is $0$, and the next bit is $1$.
\end{lemma}

\begin{definition}[Tape positions]
\lean{Complexity.LogProg.TPos}
\label{Complexity.LogProg.TPos}
\leanok
A tape position is either the left boundary, the right boundary, or a cell position
specified by a natural-number cell index and a Boolean marker.
\end{definition}

\begin{definition}[Segment word length]
\lean{Complexity.LogProg.wlen}
\label{Complexity.LogProg.wlen}
\leanok
\uses{Complexity.LogProg.Seg, Complexity.LogProg.seg}
For a list of segments and a natural-number index, the word length at that index is the
length of the Boolean list in the segment selected there, using the default segment when
the index is out of range.
\end{definition}

\begin{definition}[Validity of a tape position]
\lean{Complexity.LogProg.TPos.Valid}
\label{Complexity.LogProg.TPos.Valid}
\leanok
\uses{Complexity.LogProg.Seg, Complexity.LogProg.TPos, Complexity.LogProg.wlen}
Given a list of segments and a natural-number index, a cell position with cell index $c$
is valid exactly when $c$ is less than the word length at that index. The left and right
boundary positions are always valid.
\end{definition}

\begin{definition}[Cell offset]
\lean{Complexity.LogProg.cellOff}
\label{Complexity.LogProg.cellOff}
\leanok
\uses{Complexity.LogProg.Seg, Complexity.LogProg.seg}
Given a list of segments, an index $s$, a cell index $c$, and a Boolean marker $p$, the
cell offset is $2c+p$ when the selected segment's flag is true, and is $c$ when that
flag is false; here a Boolean is identified with $0$ or $1$.
\end{definition}

\begin{definition}[Position index]
\lean{Complexity.LogProg.vpos}
\label{Complexity.LogProg.vpos}
\leanok
\uses{Complexity.LogProg.Seg, Complexity.LogProg.TPos, Complexity.LogProg.cellOff, Complexity.LogProg.off, Complexity.LogProg.rlen, Complexity.LogProg.seg}
For a list of segments and an index $s$, assign the left boundary the offset preceding
segment $s$, and assign the right boundary that offset plus one and the selected
segment's encoded length. A cell with index $c$ and marker $p$ is assigned the preceding
offset plus one and its cell offset.
\end{definition}

\begin{lemma}[Cell Offset Bound]
\lean{Complexity.LogProg.cellOff_lt}
\label{Complexity.LogProg.cellOff_lt}
\leanok
\uses{Complexity.LogProg.Seg, Complexity.LogProg.cellOff, Complexity.LogProg.rlen, Complexity.LogProg.seg, Complexity.LogProg.wlen}
\difficulty{4}
Let a segment be a Boolean word together with a flag, and let a finite list of such
segments be given. For indices $s,c\in\mathbb{N}$ and a Boolean $p$, if $c$ is less than
the word length of the defaulted $s$th segment, then the cell offset for that segment
and position is strictly less than its rendered length.
\end{lemma}

\begin{definition}[Tape movement rule]
\lean{Complexity.LogProg.tmove}
\label{Complexity.LogProg.tmove}
\leanok
\uses{Complexity.LogProg.Seg, Complexity.LogProg.TPos, Complexity.LogProg.seg, Complexity.LogProg.wlen}
Given a list of segments, a segment index $s$, a tape position, and a direction in
$	ext{SignType}$, this rule returns a new segment index and tape position. Direction
zero leaves both unchanged. In the positive direction, at cell $c$ with marker $p$, if
the current segment's flag is true and $p$ is false, the marker becomes true; otherwise
the position advances to cell $c+1$ when that cell lies within the segment, and moves to
the right boundary when it does not. From the right boundary, it moves to the left
boundary of segment $s+1$ if that segment exists, and otherwise stays at the right
boundary; from the left boundary, it moves to cell $0$ with marker false if the segment
is nonempty, and otherwise to the right boundary. In the negative direction, at cell $c$
with marker $p$, if the current segment's flag is true and $p$ is true, the marker
becomes false; otherwise the position moves to cell $c-1$ with marker true when $c>0$,
and to the left boundary when $c=0$. From the right boundary, it moves to the last cell
with marker true if the segment is nonempty, and otherwise to the left boundary; from
the left boundary, it moves to the right boundary of segment $s-1$ when $s>0$, and
otherwise stays at the left boundary.
\end{definition}

\begin{lemma}[Validity is preserved by a tape move]
\lean{Complexity.LogProg.tmove_valid}
\label{Complexity.LogProg.tmove_valid}
\leanok
\uses{Complexity.LogProg.Seg, Complexity.LogProg.TPos, Complexity.LogProg.TPos.Valid, Complexity.LogProg.tmove}
\difficulty{4}
Let a finite sequence of segments be given, each consisting of a Boolean word and a
Boolean flag. For a segment index $s$ within this sequence, suppose a tape position is
valid: positions at the left or right boundary are always valid, and a position at a
cell is valid when its cell index is less than that segment's word length. For any move
direction, the resulting segment index remains within the sequence, and the resulting
position is valid for that segment.
\end{lemma}

\begin{lemma}[Right-boundary position of the final segment]
\lean{Complexity.LogProg.vpos_right_last}
\label{Complexity.LogProg.vpos_right_last}
\leanok
\uses{Complexity.LogProg.Seg, Complexity.LogProg.length_vword, Complexity.LogProg.vpos, Complexity.LogProg.vword}
\difficulty{3}
Let $s$ be the index of the final segment in a finite list of segments. The virtual
position immediately to the right of that segment is one position beyond the length of
the concatenated virtual word: it equals $|v|+1$, where $v$ is the virtual word formed
from the list.
\end{lemma}

\begin{lemma}[Offset and segment length bound]
\lean{Complexity.LogProg.off_add_rlen_le}
\label{Complexity.LogProg.off_add_rlen_le}
\leanok
\uses{Complexity.LogProg.Seg, Complexity.LogProg.length_render, Complexity.LogProg.off, Complexity.LogProg.rlen, Complexity.LogProg.seg, Complexity.LogProg.seg_cons_succ, Complexity.LogProg.vword}
\difficulty{4}
Let $segs$ be a finite list of segments and let $j$ be an index with $j<|segs|$. The
offset before segment $j$, plus the length of that segment after rendering, is at most
the length of the virtual word formed from the list.
\end{lemma}

\begin{lemma}[Value after moving the input head]
\lean{Complexity.LogProg.moveInputPos_val}
\label{Complexity.LogProg.moveInputPos_val}
\leanok
\uses{Turing.moveInputPos}
\difficulty{4}
For natural numbers $N,p$ with $p<N+2$, and for a movement sign $m\in\{-1,0,1\}$, the
value of the resulting input-head position is $p$ when $m=0$, $\min(p+1,N+1)$ when
$m=1$, and $p-1$ when $m=-1$.
\end{lemma}

\begin{lemma}[Virtual head movement agrees with tape movement]
\lean{Complexity.LogProg.vpos_tmove}
\label{Complexity.LogProg.vpos_tmove}
\leanok
\uses{Complexity.LogProg.Seg, Complexity.LogProg.TPos, Complexity.LogProg.TPos.Valid, Complexity.LogProg.cellOff, Complexity.LogProg.moveInputPos_val, Complexity.LogProg.off, Complexity.LogProg.off_add_rlen_le, Complexity.LogProg.off_succ, Complexity.LogProg.rlen, Complexity.LogProg.rlen_false, Complexity.LogProg.rlen_true, Complexity.LogProg.seg, Complexity.LogProg.tmove, Complexity.LogProg.vpos, Complexity.LogProg.vpos_right_last, Complexity.LogProg.vword, Complexity.LogProg.wlen, Turing.moveInputPos}
\difficulty{6}
Let a finite list of segments, a segment index $s$, and a tape position $tp$ be given,
with $s$ in range and $tp$ valid for that segment. For either head-movement direction,
if the virtual position of $tp$ is less than the virtual-word length plus two, then the
virtual position after applying the segment-tape movement equals the position obtained
by moving the head by that direction in the bounded virtual input, with movement clamped
to its boundary positions.
\end{lemma}

\begin{definition}[Symbol at a tape position]
\lean{Complexity.LogProg.vsym}
\label{Complexity.LogProg.vsym}
\leanok
\uses{Complexity.LogProg.Seg, Complexity.LogProg.TPos, Complexity.LogProg.seg}
For a list of segments and an index $s$, the symbol at the left boundary is absent when
$s=0$ and otherwise is true. At a cell with index $c$, it is the Boolean at index $c$ in
the selected segment's list, if present. At the right boundary, it is false when a next
segment exists and is absent otherwise.
\end{definition}

\begin{lemma}[Virtual input symbol at a segment position]
\lean{Complexity.LogProg.inputSymbol_vpos}
\label{Complexity.LogProg.inputSymbol_vpos}
\leanok
\uses{Complexity.LogProg.Seg, Complexity.LogProg.TPos, Complexity.LogProg.TPos.Valid, Complexity.LogProg.cellOff, Complexity.LogProg.cellOff_lt, Complexity.LogProg.getElem?_vword_seg, Complexity.LogProg.getElem?_vword_sep, Complexity.LogProg.off, Complexity.LogProg.off_succ, Complexity.LogProg.render, Complexity.LogProg.rlen, Complexity.LogProg.seg, Complexity.LogProg.vpos, Complexity.LogProg.vpos_right_last, Complexity.LogProg.vsym, Complexity.LogProg.vword, Complexity.LogProg.wlen, Turing.Cfg, Turing.Cfg.inputSymbol, Turing.FinTM.inputSymbol_at, Turing.getElem?_dbl}
\difficulty{5}
Let $\mathrm{segs}$ be a list of segments, let $s<|\mathrm{segs}|$, and let $tp$ be a
valid virtual position in segment $s$. For any configuration whose input is the virtual
word formed from these segments, if its input-head position is the virtual offset
corresponding to $tp$, then the symbol under the input head equals the virtual symbol
specified by $tp$ in segment $s$.
\end{lemma}
